\documentclass[8pt]{extarticle}

\usepackage[margin=0.30in]{geometry}
\usepackage[T1]{fontenc}
\usepackage[utf8]{inputenc}
\usepackage{amsmath,amssymb}
\usepackage{booktabs}
\usepackage{tabularx}
\usepackage{multicol}
\usepackage{enumitem}
\usepackage{titlesec}
\usepackage[protrusion=true,expansion=false]{microtype}

\setlength{\parindent}{0pt}
\setlength{\parskip}{0.5pt}
\setlength{\columnsep}{0.16in}
\setlist[itemize]{nosep,leftmargin=*}
\setlist[enumerate]{nosep,leftmargin=*}
\renewcommand{\arraystretch}{0.92}
\setlength{\abovedisplayskip}{2pt}
\setlength{\belowdisplayskip}{2pt}
\setlength{\abovedisplayshortskip}{1pt}
\setlength{\belowdisplayshortskip}{1pt}
\titleformat{\section}{\large\bfseries}{}{0pt}{}[\vspace{-3pt}\hrule\vspace{3pt}]
\titleformat{\subsection}{\normalsize\bfseries}{}{0pt}{}
\titlespacing*{\section}{0pt}{3pt}{1pt}
\titlespacing*{\subsection}{0pt}{2pt}{1pt}

\newcolumntype{Y}{>{\raggedright\arraybackslash}X}

\begin{document}

\begin{center}
    {\Large \textbf{SEG800 Image Processing Cheatsheet}}\\[-1pt]
    {\small Based on Module 7 and Module 8 study summaries}
\end{center}

\begin{multicols}{2}

\section{Concept Study Notes}

\begin{tabularx}{\linewidth}{@{}p{0.38\linewidth}Y@{}}
\toprule
\textbf{Concept} & \textbf{Reminder} \\
\midrule
Digital image processing & Processing a digital image using a digital computer. \\
Human vision: cones & Cones handle bright-light and color vision. \\
Enhancement vs restoration & Enhancement is subjective; restoration uses mathematical/probabilistic models. \\
Affine transformations & Preserve parallelism. \\
Quantization & Converts continuous brightness values into discrete levels. \\
Weber ratio & Measures brightness discrimination. \\
Digitizer & Converts analog signals into digital form. \\
Image energy source & Most digital images use the electromagnetic spectrum. \\
Point operation & Depends only on the corresponding input pixel. \\
Negative transformation & For 8-bit images: $s = 255 - z$. \\
Neighborhood operation & Image smoothing uses neighboring pixels. \\
Geometric transformation & Changes pixel locations. \\
Histogram equalization & Enhances contrast by spreading/flattening the histogram. \\
Gamma correction & Mainly changes brightness/intensity perception. \\
Translation & $u = x + t_x$, $v = y + t_y$. \\
\bottomrule
\end{tabularx}

\section{Core Definitions}

\textbf{Image:} a 2D function $f(x,y)$, where $x,y$ are spatial coordinates and $f(x,y)$ is intensity or gray level.

\textbf{Digital image:} an image where $x$, $y$, and intensity values are finite and discrete.

\textbf{Pixel:} one image element with location $(x,y)$ and intensity $f(x,y)$.

\textbf{Image enhancement:} improves an image for a specific application; subjective.

\textbf{Image restoration:} improves an image using objective mathematical/probabilistic degradation models.

\textbf{Sampling:} converts continuous spatial coordinates into discrete coordinates; related to spatial resolution.

\textbf{Quantization:} converts continuous brightness/intensity values into discrete digital values; related to intensity resolution.

\textbf{Point operation:} each output pixel depends only on the same input pixel.
\[
I_{out}(i,j)=f(I_{in}(i,j))
\]

\textbf{Neighborhood operation:} output depends on a pixel and its neighbors. Example: smoothing/blurring.

\textbf{Geometric transformation:} maps $(x,y)$ to $(u,v)$; changes locations, not intensity values.

\textbf{Affine transformation:} preserves parallelism. Examples: translation, scaling, rotation, shear.

\section{Human Vision}

\textbf{Cones:} bright-light/photopic vision; sensitive to color.\\
\textbf{Rods:} dim-light/scotopic vision; sensitive to low illumination.\\
\textbf{Blind spot:} area with optic nerve and no receptors.

\section{Formula Box}

\textbf{Weber ratio}
\[
\frac{\Delta I_c}{I}
\]
$\Delta I_c$ is the minimum detectable change in illumination; $I$ is background illumination.

\textbf{Brightness levels}
\[
2^b
\]
$b$ is number of bits. Example: $8$ bits gives $2^8=256$ levels.

\textbf{Negative transformation}
\[
s=255-z
\]
$z$ is input intensity; $s$ is output intensity.

\textbf{Gamma correction}
\[
s=cr^\gamma
\]
$r$ is input intensity, $s$ is output intensity, $c$ is a constant, and $\gamma$ is gamma.

\textbf{Lens formula}
\[
\frac{1}{z_o}+\frac{1}{z_i}=\frac{1}{f}
\]
$z_o$ is object distance, $z_i$ is image distance, and $f$ is focal length.

\textbf{Histogram}
\[
h(r_k)=\text{number of pixels with intensity } r_k
\]

\textbf{Normalized histogram}
\[
p(r_k)=\frac{h(r_k)}{MN}
\]

\textbf{Histogram equalization}
\[
s_k=T(r_k)=(L-1)\sum_{j=0}^{k}p(r_j)
\]
$L$ is the number of intensity levels.

\section{Geometric Transform Formulas}

Homogeneous coordinate form:
\[
\begin{bmatrix}u\\v\\1\end{bmatrix}
=
\begin{bmatrix}
h_{11}&h_{12}&h_{13}\\
h_{21}&h_{22}&h_{23}\\
h_{31}&h_{32}&h_{33}
\end{bmatrix}
\begin{bmatrix}x\\y\\1\end{bmatrix}
\]

\textbf{Translation}
\[
u=x+t_x,\qquad v=y+t_y
\]
\[
\begin{bmatrix}1&0&t_x\\0&1&t_y\\0&0&1\end{bmatrix}
\]

\textbf{Scaling}
\[
u=s_xx,\qquad v=s_yy
\]
\[
\begin{bmatrix}s_x&0&0\\0&s_y&0\\0&0&1\end{bmatrix}
\]

\textbf{Rotation, counterclockwise}
\[
u=x\cos\theta-y\sin\theta,\qquad v=x\sin\theta+y\cos\theta
\]
\[
\begin{bmatrix}\cos\theta&-\sin\theta&0\\\sin\theta&\cos\theta&0\\0&0&1\end{bmatrix}
\]

\textbf{Shear}
\[
u=x+y\,sh_x,\qquad v=y+x\,sh_y
\]
\[
\begin{bmatrix}1&sh_x&0\\sh_y&1&0\\0&0&1\end{bmatrix}
\]

\section{Pixel Adjacency}

For $p=(x,y)$:
\[
N_4(p)=\{(x+1,y),(x-1,y),(x,y+1),(x,y-1)\}
\]
\[
N_D(p)=\{(x+1,y+1),(x+1,y-1),(x-1,y+1),(x-1,y-1)\}
\]
\[
N_8(p)=N_4(p)\cup N_D(p)
\]

Two pixels $p$ and $q$ are \textbf{m-adjacent} if either:
\[
q\in N_4(p)
\]
or $q\in N_D(p)$ and $N_4(p)\cap N_4(q)$ has no pixels with values from $V$.

Checklist: confirm both pixels are in foreground set $V$; check 4-neighbor first; for diagonal pixels, inspect common 4-neighbors.

\section{Optics / Lens Formula}

\[
\frac{1}{z_o}+\frac{1}{z_i}=\frac{1}{f}
\]

Solving for object distance:
\[
\frac{1}{z_o}=\frac{1}{f}-\frac{1}{z_i}
\]
\[
z_o=\frac{1}{\frac{1}{f}-\frac{1}{z_i}}
\]

Keep all distances in the same unit before calculating. Convert at the end if needed.

\section{Homography Composition}

If applying scale $\rightarrow$ rotate $\rightarrow$ translate, multiply in reverse application order:
\[
H=TRS
\]

Apply using:
\[
\begin{bmatrix}u\\v\\1\end{bmatrix}=H\begin{bmatrix}x\\y\\1\end{bmatrix}
\]

Geometric transformations move pixel locations but do not directly change pixel intensity values. Raster images may need interpolation or rounding.

\section{Extra Study Topics}

\textbf{Processing levels}

\begin{tabularx}{\linewidth}{@{}p{0.28\linewidth}p{0.16\linewidth}Y Y@{}}
\toprule
\textbf{Level} & \textbf{Input} & \textbf{Output} & \textbf{Examples} \\
\midrule
Image processing & Image & Image & Noise reduction, contrast enhancement, sharpening \\
Image analysis & Image & Features & Edges, contours, objects \\
Computer vision & Image & Meaning & Recognition, scene understanding \\
\bottomrule
\end{tabularx}

\textbf{Image acquisition pipeline}
\begin{align*}
&\text{Energy source}\rightarrow\text{Scene/object}\rightarrow\text{Imaging system}\\
&\rightarrow\text{Image plane}\rightarrow\text{Digitized image}
\end{align*}

Most images use the electromagnetic spectrum: visible light, X-ray, infrared, ultraviolet, microwave, and radio waves. Other sources include acoustic waves, ultrasonic waves, electron beams, and synthetic images.

\textbf{Digital image types}

\begin{tabularx}{\linewidth}{@{}p{0.32\linewidth}Y@{}}
\toprule
\textbf{Type} & \textbf{Meaning} \\
\midrule
Binary & 2D array of 0s and 1s \\
Grayscale & 2D array of brightness values, often 0--255 or normalized 0--1 \\
Color & Usually three 2D arrays, such as R, G, B \\
Multispectral & Includes frequencies beyond visible light \\
Depth map & Stores distance/depth information \\
\bottomrule
\end{tabularx}

\textbf{Brightness adaptation:} the eye changes sensitivity depending on current light level. It adapts to a large intensity range, but not all at once.

\textbf{Weber ratio interpretation}
\[
\text{Large Weber ratio}\rightarrow\text{poor brightness discrimination}
\]
\[
\text{Small Weber ratio}\rightarrow\text{better brightness discrimination}
\]

\textbf{Visual perception effects:} Mach band effect, simultaneous contrast, optical illusions.

\textbf{Saturation:} highest intensity value beyond which values are clipped.\\
\textbf{Noise:} random unwanted variation in image intensity.

\textbf{Region terms:} a path is a sequence of connected pixels; a connected set/region has paths between pixels inside it; a boundary/contour contains pixels adjacent to outside pixels.

\textbf{Distance measures} for $p=(x,y)$ and $q=(u,v)$:
\[
D_e(p,q)=\sqrt{(x-u)^2+(y-v)^2}
\]
\[
D_4(p,q)=|x-u|+|y-v|
\]
\[
D_8(p,q)=\max(|x-u|,|y-v|)
\]

\textbf{Linear operation property}
\[
H[af_1(x,y)+bf_2(x,y)]=aH[f_1(x,y)]+bH[f_2(x,y)]
\]

\textbf{Image arithmetic}
\[
s=f+g,\quad d=f-g,\quad p=f\times g,\quad v=\frac{f}{g}
\]

\textbf{Linear blend / weighted addition}
\[
\text{dst}=\alpha\,\text{img}_1+\beta\,\text{img}_2+\gamma
\]
$\alpha$ and $\beta$ are image weights; $\gamma$ is a brightness offset.

\textbf{Brightness and contrast point operations}
\[
I_{out}(i,j)=I_{in}(i,j)+b
\]
\[
I_{out}(i,j)=aI_{in}(i,j)
\]
Adding changes brightness; multiplying changes contrast.

\textbf{Contrast stretching:} expands intensity values to improve contrast.

\textbf{Thresholding:} separates pixels into dark/light groups.
\[
\text{pixel}\ge\text{threshold}\rightarrow\text{white},\qquad
\text{pixel}<\text{threshold}\rightarrow\text{black}
\]

\textbf{Gamma rule}
\[
\gamma<1\rightarrow\text{brighter},\qquad \gamma>1\rightarrow\text{darker}
\]

\textbf{Transformation families}

\begin{tabularx}{\linewidth}{@{}p{0.34\linewidth}Y@{}}
\toprule
\textbf{Transformation} & \textbf{Preserves} \\
\midrule
Translation & Orientation \\
Rigid / Euclidean & Lengths \\
Similarity & Angles \\
Affine & Parallelism \\
Projective & Straight lines; parallel lines may not stay parallel \\
\bottomrule
\end{tabularx}

\textbf{Interpolation:} assigns intensity when transformed locations are non-integer.

\begin{tabularx}{\linewidth}{@{}p{0.34\linewidth}Y@{}}
\toprule
\textbf{Type} & \textbf{Meaning} \\
\midrule
Nearest-neighbor & Closest pixel; fast but blocky \\
Bilinear & Nearby pixels; smoother \\
Bicubic & More neighbors; smoother but slower \\
\bottomrule
\end{tabularx}

\section{Quick Memory Tips}

\begin{itemize}
    \item Enhancement = subjective; restoration = model-based.
    \item Cones = color/bright light; rods = dim light.
    \item Sampling = spatial coordinates; quantization = intensity values.
    \item Point operation = same pixel only; neighborhood operation = pixel plus neighbors.
    \item Geometric transformation = changes locations; pixel value stays the same.
    \item Affine preserves parallelism; projective preserves straight lines.
    \item Histogram equalization = contrast enhancement.
    \item Gamma correction = brightness perception; $\gamma<1$ brighter, $\gamma>1$ darker.
    \item Negative image: $s=255-z$.
    \item Translation: add $t_x$ and $t_y$.
    \item Large Weber ratio = poor brightness discrimination; small Weber ratio = better discrimination.
\end{itemize}

\end{multicols}

\end{document}
